\documentclass[12pt]{article}
\usepackage[margin=0.58in,top=0.62in,bottom=0.48in]{geometry}
\usepackage{lmodern}
\usepackage{microtype}
\usepackage{fancyhdr}
\usepackage{titlesec}
\usepackage{enumitem}
\usepackage{lastpage}
\usepackage[hidelinks]{hyperref}

\setlength{\parindent}{0pt}
\setlength{\parskip}{0.18em}
\setlength{\headheight}{26pt}
\raggedbottom
\titleformat{\section}{\normalsize\bfseries\scshape}{}{0pt}{}
\titleformat{\subsection}{\normalsize\bfseries}{}{0pt}{}
\titlespacing*{\section}{0pt}{0.35em}{0.12em}
\titlespacing*{\subsection}{0pt}{0.25em}{0.08em}
\pagestyle{fancy}
\fancyhf{}
\fancyhead[C]{\footnotesize Lit Example Reader \quad Assignment ID: U1L17LE \quad Page \thepage\ of \pageref{LastPage}}
\renewcommand{\headrulewidth}{0.5pt}

\newenvironment{playpassage}{\begin{quote}\footnotesize\setlength{\parskip}{0.08em}\setlength{\parindent}{0pt}}{\end{quote}}
\newcommand{\speaker}[1]{\textbf{#1.}}

\begin{document}

\begin{center}
{\LARGE\bfseries Week 17 Lit Example Reader}\par\vspace{0.02em}
{\large\scshape Shakespeare: Reasoning About Oneself}\par\vspace{0.06em}
\rule{0.78\textwidth}{0.5pt}
\end{center}

\textbf{Purpose:} Shakespeare passages for Week 17: apply a multi-tool self-deception autopsy to characters who justify themselves.

\textbf{What to watch for:} Desire, private conclusions, cleaner public reasons, shifted terms, missing proof, and rhetoric used on the self.

\section*{Before the Passages: The Week 17 Logic Tool}

\begin{itemize}[leftmargin=1.3em,itemsep=0.05em,topsep=0.08em,parsep=0.03em]
\item \textbf{Self-deception autopsy:} desire; conclusion; stated reasons; hidden assumption; exposing tool; remaining need; repaired claim.
\item \textbf{Private argument / public cover:} inner justification versus cleaner outer story.
\item \textbf{Tool stack:} terms, assumptions, circularity, evidence, sources, rhetoric-vs-proof, and more as needed.
\end{itemize}

\section*{Passage 1: Macbeth, Act 1, Scene 7 --- Macbeth Before the Murder}

\textbf{Context:} Alone, Macbeth weighs the murder of Duncan. He sees moral reasons against the act, then feels vaulting ambition. The speech shows a mind arguing with itself before Lady Macbeth returns.

\begin{playpassage}
\speaker{Macbeth} If it were done when 'tis done, then 'twere well\\
It were done quickly: if the assassination\\
Could trammel up the consequence, and catch\\
With his surcease success; that but this blow\\
Might be the be-all and the end-all here,\\
...\\
He's here in double trust;\\
First, as I am his kinsman and his subject,\\
Strong both against the deed; then, as his host,\\
Who should against his murderer shut the door,\\
Not bear the knife myself.\\
...\\
I have no spur\\
To prick the sides of my intent, but only\\
Vaulting ambition, which o'erleaps itself\\
And falls on the other.
\end{playpassage}

\subsection*{Reader Notes}

This is your assisted example. Macbeth can see duties that count against the murder: kinship, subject-loyalty, and host responsibility. He also names the spur: ambition. A self-deception autopsy does not deny his clarity; it asks how later choices treat what he already knows, and how desire recruits new reasons, manhood pressure, or prophecy security later. For this passage, map what he knows versus what he wants.

\textbf{Reading task:} Underline duties against the deed. Box the named spur. In the margin, write Macbeth's desire in five words or fewer.

\newpage

\section*{Passage 2: Julius Caesar, Act 2, Scene 1 --- Brutus in the Orchard}

\textbf{Context:} Brutus reasons privately that Caesar must be stopped before he becomes a serpent's egg of tyranny. He has not yet seen clear present proof of the crown-turned-tyrant he fears.

\begin{playpassage}
\speaker{Brutus} It must be by his death: and for my part,\\
I know no personal cause to spurn at him,\\
But for the general. He would be crown'd:\\
How that might change his nature, there's the question.\\
...\\
And therefore think him as a serpent's egg\\
Which, hatch'd, would, as his kind, grow mischievous,\\
And kill him in the shell.
\end{playpassage}

\subsection*{Reader Notes}

This passage shifts more work to you. Watch the movement from ``it must be by his death'' to the analogy that makes killing feel preventive and responsible. Do not retell the plot. Ask what Brutus desires for Rome and for his own honor, what is assumed rather than shown, and which earlier tools (analogy, hidden assumption, missing proof, public cover of necessity) belong in the autopsy.

\textbf{Reading task:} Underline the conclusion. Star the analogy. Write one question about what evidence would still be needed before death is required.

\section*{How to Use These Passages on the Worksheet}

Do not summarize the scenes. Use Week 17's self-deception autopsy: desire, conclusion, stated reasons, hidden assumption, exposing tool, remaining need, and repaired narrower claim. Passage 1 is the assisted model; Passage 2 requires more independent transfer.

\end{document}
